\section{Classical discrete probability laws as mechanisms}

Classical discrete distributions should not be learned as a table of formulas.
Each one is a compact description of a random mechanism. The formula for the
probability mass function is important, but it should be read as the consequence
of how the experiment is constructed.

In this section we collect several basic discrete laws and emphasize the
mechanism behind each one.

\subsection*{Bernoulli law: one success-or-failure trial}

The simplest discrete random variable records whether one event occurred. Let a
single trial have two possible outcomes: success and failure. Define
\[
X=
\begin{cases}
1, & \text{success occurs},\\
0, & \text{failure occurs}.
\end{cases}
\]
If
\[
P(\text{success})=p,
\]
then
\[
P(X=1)=p,\qquad P(X=0)=1-p.
\]

\begin{definitionbox}
A random variable \(X\) has a \textbf{Bernoulli distribution} with parameter
\(p\), written \(X\sim\operatorname{Bernoulli}(p)\), if
\[
P(X=1)=p,\qquad P(X=0)=1-p.
\]
\end{definitionbox}

The Bernoulli law is the probability law of an indicator variable. If \(X=I_A\),
then \(X\sim\operatorname{Bernoulli}(P(A))\).

Its expectation and variance are
\[
\E[X]=p,\qquad \Var(X)=p(1-p).
\]

\subsection*{Binomial law: number of successes in fixed independent trials}

Now repeat a Bernoulli trial \(n\) times. Assume that each trial has success
probability \(p\), and that the trials are independent. Let
\[
X=\text{the number of successes in }n\text{ trials}.
\]
Then
\[
\operatorname{supp}(X)=\{0,1,2,\ldots,n\}.
\]

To have \(X=k\), exactly \(k\) of the \(n\) positions must be successes. There
are
\[
\binom nk
\]
ways to choose those positions. Each sequence with \(k\) successes and \(n-k\)
failures has probability
\[
p^k(1-p)^{n-k}.
\]
Therefore
\[
P(X=k)=\binom nk p^k(1-p)^{n-k},\qquad k=0,1,\ldots,n.
\]

\begin{definitionbox}
A random variable \(X\) has a \textbf{binomial distribution} with parameters
\(n\) and \(p\), written \(X\sim\operatorname{Binomial}(n,p)\), if
\[
P(X=k)=\binom nk p^k(1-p)^{n-k},\qquad k=0,1,\ldots,n.
\]
\end{definitionbox}

The mechanism matters. Binomial means fixed number of trials, two outcomes per
trial, constant success probability, independence, and counting successes.

Using indicators, if
\[
X=I_1+I_2+\cdots+I_n,
\]
where each \(I_i\sim\operatorname{Bernoulli}(p)\), then
\[
\E[X]=np.
\]
With independence one also obtains
\[
\Var(X)=np(1-p).
\]

\subsection*{Geometric law: waiting for the first success}

In a binomial model, the number of trials is fixed and the number of successes is
random. In the geometric model, the number of successes is fixed at one first
success, and the waiting time is random.

Repeat independent Bernoulli trials with success probability \(p\). Let
\[
N=\text{the trial on which the first success occurs}.
\]
Then
\[
\operatorname{supp}(N)=\{1,2,3,\ldots\}.
\]
For \(N=k\), the first \(k-1\) trials must be failures and the \(k\)-th trial
must be a success. Hence
\[
P(N=k)=(1-p)^{k-1}p,\qquad k=1,2,3,\ldots.
\]

\begin{definitionbox}
A random variable \(N\) has a \textbf{geometric distribution} with parameter
\(p\), in the waiting-time convention, if
\[
P(N=k)=(1-p)^{k-1}p,\qquad k=1,2,3,\ldots.
\]
\end{definitionbox}

The geometric distribution is a model for waiting until the first success. Its
support is infinite because success may be delayed for any number of trials.

\subsection*{Hypergeometric law: sampling without replacement}

Suppose a population contains \(N\) objects. Among them, \(K\) are classified as
successes and \(N-K\) as failures. We sample \(n\) objects without replacement and
ignore order. Let
\[
X=\text{the number of successes in the sample}.
\]

To get \(X=k\), choose \(k\) successes from the \(K\) success objects and choose
\(n-k\) failures from the \(N-K\) failure objects. The total number of samples is
\(\binom Nn\). Therefore
\[
P(X=k)=\frac{\binom Kk\binom{N-K}{n-k}}{\binom Nn}.
\]

\begin{definitionbox}
A random variable \(X\) has a \textbf{hypergeometric distribution} with
parameters \(N,K,n\) if
\[
P(X=k)=\frac{\binom Kk\binom{N-K}{n-k}}{\binom Nn}
\]
for the admissible values of \(k\).
\end{definitionbox}

This model differs from the binomial model because sampling is done without
replacement. The draws are dependent: after one object is selected, the
composition of the remaining population changes.

\subsection*{Poisson law: counts of rare events}

The Poisson distribution is used to model counts of events occurring in a fixed
region of time, space, area, or volume, when events occur with a stable average
rate and without strong clustering beyond what the model allows.

Let
\[
X=\text{the number of events in one observation window}.
\]
The support is
\[
\{0,1,2,\ldots\}.
\]
A Poisson model with parameter \(\lambda>0\) has mass function
\[
P(X=k)=e^{-\lambda}\frac{\lambda^k}{k!},\qquad k=0,1,2,\ldots.
\]

\begin{definitionbox}
A random variable \(X\) has a \textbf{Poisson distribution} with parameter
\(\lambda>0\), written \(X\sim\operatorname{Poisson}(\lambda)\), if
\[
P(X=k)=e^{-\lambda}\frac{\lambda^k}{k!},\qquad k=0,1,2,\ldots.
\]
\end{definitionbox}

The parameter \(\lambda\) is both the expectation and the variance:
\[
\E[X]=\lambda,\qquad \Var(X)=\lambda.
\]
In this chapter, the main point is interpretation: the Poisson law models a
count, not a waiting time and not a fixed number of trials.

\subsection*{Comparing the mechanisms}

The following table summarizes the roles of the classical laws.
\[
\begin{array}{c|c|c}
\text{Law} & \text{Random variable} & \text{Mechanism}\\
\hline
\text{Bernoulli} & 0\text{ or }1 & \text{one success-or-failure trial}\\
\text{Binomial} & \text{number of successes} & \text{fixed independent trials}\\
\text{Geometric} & \text{waiting time} & \text{repeat until first success}\\
\text{Hypergeometric} & \text{number of successes} & \text{sampling without replacement}\\
\text{Poisson} & \text{count of events} & \text{events in a fixed window}
\end{array}
\]

The same numerical value may appear in several models, but the interpretation is
different. A value \(k\) in a binomial model means \(k\) successes among a fixed
number of trials. A value \(k\) in a geometric model means the first success
occurred on trial \(k\). A value \(k\) in a Poisson model means \(k\) events were
counted in a window.

\begin{remarkbox}
Choosing a distribution is not choosing a formula. It is choosing a mechanism.
Before using a classical law, identify what is being counted, what is fixed, what
is random, and what assumptions connect the experiment to the model.
\end{remarkbox}
